\section{Problem Summary}
Given a string that may contain repeated characters, list every distinct permutation in lexicographic order and print how many such permutations there are.

\section{Editorial}
The duplicates are the only complication. If every character were distinct, ordinary permutation generation would be enough. With repeated characters, we must avoid printing the same arrangement more than once.

Sorting the string first solves both requirements at once:
\begin{itemize}
\item the first arrangement becomes the lexicographically smallest one
\item repeated calls to \texttt{next\_permutation} visit each distinct arrangement in sorted order
\end{itemize}

This works because equal characters are indistinguishable. Starting from the sorted string ensures the algorithm walks through the multiset permutations without manual deduplication tables.

Since the input is small, storing all permutations before printing them is completely reasonable and makes it easy to print the count first.

\section{Pseudocode}
Sort the string.

Create an empty list of answers.

Do:
\begin{itemize}
\item add the current string to the answer list
\end{itemize}
while another lexicographically larger permutation exists.

Print the number of stored strings, then print each one.

\section{Complexity Analysis}
\begin{itemize}
\item Time: \(O(k \cdot n)\), where \(k\) is the number of distinct permutations
\item Memory: \(O(k \cdot n)\) when storing the permutations before printing
\end{itemize}

\section{Notes}
For this problem size, the output dominates everything else, so a simple library-based permutation loop is the right tradeoff.
